\scanpage{12}
(en plus des deux diagonales et des deux codiagonales elles-mêmes, qui font encore quatre autres ensembles à deux éléments, mais groupés par paires ``indistinguables'' \ldots). On a des bijections canoniques
\[
\underline\omega_S\simeq C,\qquad\underline\omega_A\simeq D
\]
provenant du fait qu'on a
\[
\underline\omega_S\simeq\underline\omega_{V_S^*}=\underline\omega_{\{\underline a\}\amalg C}\simeq\underline\omega_C\simeq C
\]
La convention pour orienter un ensemble à trois éléments $J$ (ici $V_S^*$) dont un élément $a$ (ici $\underline a$) est fixé, en termes du choix d'un élément \inserted{$x$} du complémentaire (ici $V_S^*\setminus\{a\}=C$) est de prendre l'orientation $(a,x,y)$ \deleted{(\unsure{où} $\unsure{y},\{x,y\}=V_S^{\unsure{a|*}}\setminus\{a,$} où $\{y\}=J\setminus\{a,x\}$.

On a une relation remarquable
\[
\underline\omega_Q\wedge C\wedge D\simeq\mathbb{1}_{\mathbb{F}_2}\tag*{($*$)}
\]
\margin{i.e. une restriction du groupe structural $\mathbb{F}_2^3$ de $\underline\omega_{\unsure{S|Q}}\times C\times D$ au sous-groupe $V_{\underline3}(\mathbb{F}_2)$ des $x,y,z$ tels que $x+y+z=0$}
qui provient du fait que $\underline\omega_{\unsure{S|Q}}, C,D$ sont les images inverses, dans l'extension $\unsure{\mathfrak{D}|\mathscr{D}}$ de $\mathfrak{D}/\{1,\underline a\}=\mathfrak{D}_{ab}$ par $\{1,\underline a\}$, des trois éléments non nuls de $\mathfrak{D}_{ab}$, dont la \deleted{produit} \inserted{somme} est nulle. \inserted{Pour} \deleted{\unsure{d}}es éléments $\varpi,c,d$ de $\underline\omega_Q\times C\times D$ \deleted{\unsure{satisf}} on a
\[
\varpi\wedge c\wedge d=0\quad\Longleftrightarrow\quad\varpi c d=1\text{ dans }\mathfrak{D}\tag*{($\unsure{*|**}$)}
\]
\margin{[Attention, cette relation \unsure{serait|devrait} inversée si on faisait une transposition sur l'ordre des facteurs -- il y a donc \deleted{\unsure{dans}} dans \deleted{\unsure{cette}} le choix de l'iso. $(*)$ une question de convention.}

\margin{Si on passe du carré $Q$ au carré dual $Q'$, le choix de l'iso \deleted{\unsure{$\omega$}} $\underline\omega_{Q'}\deleted{\unsure{\wedge}} C'\wedge D'\deleted{\unsure{\simeq1}}(\simeq\underline\omega_Q\wedge D\wedge C)\simeq\mathbb{1}$ est \uline{opposé} de celui pour $Q$\ldots}
% archive sheet number: 11
